\documentclass[aspectratio=169]{beamer}
\usetheme{Madrid}
\usecolortheme{beaver}
\usepackage[utf8]{inputenc}
\usepackage[T1]{fontenc}
\usepackage{tikz}
\usepackage{amssymb}
\usepackage{graphicx}
\setbeamertemplate{navigation symbols}{}
\setbeamercolor{title}{fg=white,bg=cyan!60!black}
\setbeamercolor{frametitle}{fg=white,bg=cyan!50!black}
\setbeamercolor{block title}{bg=cyan!15,fg=black}
\setbeamercolor{block body}{bg=cyan!5,fg=black}

\title{G1 — Sprite Start}
\subtitle{Coordinates, sequences — draw and move your first sprite}
\author{ByteForge --- Grade 7 --- Game Maker --- Week Overview}
\date{}

\begin{document}

\begin{frame}
    \titlepage
\end{frame}

\begin{frame}{This Week's Plan}
    \begin{itemize}
        \item Lesson 1: First Sprite
        \item Lesson 2: Four Corners
        \item Lesson 3: Rectangle Room
        \item Lesson 4: Three in a Row
        \item[$\bigstar$] BOSS --- Sprite Follows You (\texttt{lessonBoss.tex})
    \end{itemize}
\end{frame}

\begin{frame}{Real-World Example}
    \Large
    Your phone's touchscreen reporting exactly WHERE you tapped (not just THAT you tapped) — every touch event carries an (x, y), same idea as this chapter's sprite position.
\end{frame}


\begin{frame}{Let Them Discover (1/3)}
    \Large
    On paper, draw a 480x260 rectangle. Mark where (0,0), (240,130), and (479,259) actually land. Which corner is (0,0) — and why might that surprise someone used to maths graphs?
\end{frame}


\begin{frame}{Let Them Discover (2/3)}
    \Large
    How do you actually choose WHICH corner is "the" corner for a rectangle? Try this before opening the platform: agree as a class that (x,y) always means the TOP-LEFT corner — then everyone should get the same answer for where a 100x50 box at (200,80) ends.
\end{frame}


\begin{frame}{Let Them Discover (3/3)}
    \Large
    Live demo — a moving cursor readout: open grade7-coord-cursor.html (in this same lessons folder) and project it. Move your mouse around the grid and read the live (x,y) label out loud as you go — this answers "how do I know where I even am" before anyone touches real code.
\end{frame}

\begin{frame}{Files in This Folder}
    \small
    \begin{itemize}
        \item \texttt{lesson1.tex} --- First Sprite
        \item \texttt{lesson2.tex} --- Four Corners
        \item \texttt{lesson3.tex} --- Rectangle Room
        \item \texttt{lesson4.tex} --- Three in a Row
        \item \texttt{lessonBoss.tex} --- BOSS: Sprite Follows You
    \end{itemize}
    \vspace{0.5em}
    \normalsize Each lesson file is its own short deck --- open the one for
    the specific day you're teaching, not this whole overview.
\end{frame}

\end{document}
